\chapter{直角三角形中的三角函数}\label{ch:g9:trig}

在\omterm{def:g6:shapes:triangles}{直角三角形}中，只要知道一个锐角，三角形的\emph{形状}便已确定
——只有大小可以变化。三角形各边之比
因此只取决于该角：这些比分别是\omterm{def:g9:trig:ratios}{余弦}、\omterm{def:g9:trig:ratios}{正弦}和
\omterm{def:g9:trig:ratios}{正切}。结合勾股定理，
只需极少的已知信息，即可求出\omterm{def:g6:shapes:triangles}{直角三角形}的
每一条边和每一个角。

\section{再论勾股定理}

\begin{theorem}[毕达哥拉斯]\label{thm:g9:trig:pythagoras}
在\omterm{def:g6:shapes:triangles}{直角三角形} $ABC$ 中（$A$ 为\omterm{def:g3:shapes:rightangle}{直角}），斜边的平方
等于两条直角边的平方和：
\[
BC^2 = AB^2 + AC^2 .
\]
反之，若三角形中 $BC^2 = AB^2 + AC^2$，则该三角形
以 $A$ 为\omterm{def:g3:shapes:rightangle}{直角}。
\end{theorem}
\admitted

\begin{example}\label{ex:g9:trig:pythagoras}
一个\omterm{def:g6:shapes:triangles}{直角三角形}的两条直角边为 $AB = 5$ 与 $AC = 12$。于是
\[
BC^2 = 5^2 + 12^2 = 25 + 144 = 169,
\qquad\text{故}\quad BC = \sqrt{169} = 13 .
\]
求一条\emph{直角边}：若 $BC = 10$，$AB = 6$，则
$AC^2 = BC^2 - AB^2 = 100 - 36 = 64$，故 $AC = 8$。
\end{example}

\section{余弦、正弦、正切}

在\omterm{def:g6:shapes:triangles}{直角三角形}中，对于锐角 $\theta$，三条边
各有名称：\emph{斜边}（与\omterm{def:g3:shapes:rightangle}{直角}相对的边，最长
的边）、$\theta$ 的\emph{对边}，以及
$\theta$ 的\emph{邻边}（与 $\theta$ 相接的直角边）。

\begin{definition}[锐角三角函数]\label{def:g9:trig:ratios}
对\omterm{def:g6:shapes:triangles}{直角三角形}的锐角 $\theta$，规定：
\[
\cos\theta = \frac{\text{邻边}}{\text{斜边}},
\qquad
\sin\theta = \frac{\text{对边}}{\text{斜边}},
\qquad
\tan\theta = \frac{\text{对边}}{\text{邻边}} .
\]
这些比值\index{余弦}\index{正弦}\index{正切}仅取决于
$\theta$，而与三角形的大小无关（具有相同锐角的\omterm{def:g6:shapes:triangles}{直角三角形}
彼此是按比例缩放的，由泰勒斯定理可知）。
\end{definition}

\begin{omfigure}
\begin{tikzpicture}[scale=1.15]
  \coordinate (A) at (0,0);
  \coordinate (B) at (4.2,0);
  \coordinate (C) at (4.2,2.2);
  \draw[thick] (A) -- (B) -- (C) -- cycle;
  \draw[thin] (B) ++(-0.28,0) -- ++(0,0.28) -- ++(0.28,0);
  \draw[omThm, ->] (0.9,0) arc (0:27.7:0.9);
  \node[omThm, font=\small] at (1.25,0.21) {$\theta$};
  \node[below, font=\small, omDef] at (2.1,0) {邻边};
  \node[right, font=\small, omMeth] at (4.2,1.1) {对边};
  \node[above left=-2pt, font=\small, omProp] at (2.2,1.25) {斜边};
  \node[below left, font=\small] at (A) {$B$};
  \node[below right, font=\small] at (B) {$A$};
  \node[above right, font=\small] at (C) {$C$};
\end{tikzpicture}

{\small 从角 $\theta = \widehat B$ 出发看三条边：
斜边 $[BC]$、邻边 $[BA]$、对边 $[AC]$。}
\end{omfigure}

\begin{proposition}[基本性质]\label{prop:g9:trig:properties}
对每个锐角 $\theta$，有：
\begin{enumerate}
  \item $0 < \cos\theta < 1$ 且 $0 < \sin\theta < 1$（直角边短于
        斜边）；
  \item $\cos^2\theta + \sin^2\theta = 1$；
  \item $\tan\theta = \dfrac{\sin\theta}{\cos\theta}$。
\end{enumerate}
\end{proposition}

\begin{proof}
设 $a$ 为邻边，$o$ 为对边，$h$ 为
角为 $\theta$ 的\omterm{def:g6:shapes:triangles}{直角三角形}的斜边。

1. 由 $0 < a < h$ 与 $0 < o < h$，两个比值都严格介于
$0$ 与 $1$ 之间。

2. 由毕达哥拉斯定理，$a^2 + o^2 = h^2$。两边同除以 $h^2$：
\[
\cos^2\theta + \sin^2\theta
= \frac{a^2}{h^2} + \frac{o^2}{h^2}
= \frac{a^2 + o^2}{h^2} = 1 .
\]

3. $\dfrac{\sin\theta}{\cos\theta} = \dfrac{o/h}{a/h} = \dfrac oa
= \tan\theta$.
\end{proof}

\begin{method}[求边长]\label{met:g9:trig:side}
已知一条边和一个锐角，求未知边时，按以下步骤进行：
\begin{enumerate}
  \item 标出已知角，并标注三条边（斜边、
        对边、邻边）\emph{从该角出发}；
  \item 选取涉及已知边和所求边的比值
        （cos、sin 或 tan）；
  \item 列出方程并求解；
  \item 检验结果：斜边应最长。
\end{enumerate}
\end{method}

\begin{example}\label{ex:g9:trig:side}
在\omterm{def:g6:shapes:triangles}{直角三角形}中，斜边长为 $8$，其中一个角为
$35^\circ$，求该角的对边。涉及
对边与斜边的比值为\omterm{def:g9:trig:ratios}{正弦}：
\[
\sin 35^\circ = \frac{\text{对边}}{8}
\quad\Longrightarrow\quad
\text{对边} = 8 \sin 35^\circ \approx 8 \times 0.574 \approx 4.6 .
\]
对于邻边，使用\omterm{def:g9:trig:ratios}{余弦}：
$8\cos 35^\circ \approx 8 \times 0.819 \approx 6.6$。可用
毕达哥拉斯定理检验：$4.6^2 + 6.6^2 \approx 21 + 43 \approx 64 = 8^2$。
\end{example}

\begin{method}[求角]\label{met:g9:trig:angle}
已知两边时，求出它们所成的比值，
将其判定为未知角的\omterm{def:g9:trig:ratios}{余弦}、\omterm{def:g9:trig:ratios}{正弦}或\omterm{def:g9:trig:ratios}{正切}，
并对该比值使用计算器的反函数（$\cos^{-1}$、$\sin^{-1}$ 或 $\tan^{-1}$）。
\end{method}

\begin{example}\label{ex:g9:trig:angle}
一架长 $5$ m 的梯子斜靠在墙上，梯脚离墙
的距离为 $1.4$ m。梯子与地面所成的角 $\theta$ 满足
\[
\cos\theta = \frac{1.4}{5} = 0.28,
\qquad\text{故}\quad
\theta = \cos^{-1}(0.28) \approx 73.7^\circ .
\]
\end{example}

\begin{omfigure}
\begin{tikzpicture}[scale=0.85]
  \draw[thick] (0,0) -- (1.4,0);
  \draw[thick] (1.4,0) -- (1.4,4.0);
  \draw[omDef, very thick] (0,0) -- (1.4,4.0);
  \draw[thin] (1.4,0) ++(-0.25,0) -- ++(0,0.25) -- ++(0.25,0);
  \draw[omThm, ->] (0.75,0) arc (0:70.7:0.75);
  \node[omThm, font=\small] at (1.15,0.55) {$\theta$};
  \node[below, font=\small] at (0.7,0) {$1.4$ m};
  \node[omDef, font=\small, left] at (0.62,2.1) {$5$ m};
  \draw[gray, thin] (-0.7,0) -- (2.3,0);
\end{tikzpicture}

{\small 梯子问题：邻边（$1.4$ m）与
斜边（$5$ m）已知，故由\omterm{def:g9:trig:ratios}{余弦}即可求出该角。}
\end{omfigure}

\begin{example}[特殊角]\label{ex:g9:trig:special}
两个三角形给出特殊角的精确值。边长为 1 的正方形的\omterm{def:g3:division:half}{一半}
（\omterm{def:g3:shapes:rightangle}{直角}\omterm{def:g6:shapes:triangles}{等腰三角形}）斜边长为 $\sqrt2$，锐角均为 $45^\circ$：
\[
\cos 45^\circ = \sin 45^\circ = \frac{1}{\sqrt2} = \frac{\sqrt2}{2},
\qquad \tan 45^\circ = 1 .
\]
边长为 1 的\omterm{def:g6:shapes:triangles}{等边三角形}的\omterm{def:g3:division:half}{一半}给出角 $30^\circ$ 与
$60^\circ$：
\[
\cos 60^\circ = \sin 30^\circ = \frac12,
\qquad
\cos 30^\circ = \sin 60^\circ = \frac{\sqrt3}{2}.
\]
\end{example}

\section{练习}

\begin{exercise}[$\star$]\label{exo:g9:trig:1}
一个\omterm{def:g6:shapes:triangles}{直角三角形}的两直角边长分别为 $9$ 和 $12$，求斜边长。另一个
\omterm{def:g6:shapes:triangles}{直角三角形}斜边长为 $17$，一条直角边为 $8$，求另一条直角边。
\end{exercise}

\begin{exercise}[$\star$]\label{exo:g9:trig:2}
一个三角形的三边长为 $7$、$24$ 和 $25$，判断它是否为\omterm{def:g6:shapes:triangles}{直角三角形}。同样
对三边长为 $5$、$6$ 和 $8$ 的三角形作判断。
\end{exercise}

\begin{exercise}[$\star$]\label{exo:g9:trig:3}
在\omterm{def:g6:shapes:triangles}{直角三角形} $DEF$ 中，$D$ 为\omterm{def:g3:shapes:rightangle}{直角}，记 $E$ 处的角为
$\theta$，哪条边是斜边？哪条边与
$\theta$ 相对？将 $\cos\theta$、$\sin\theta$ 和 $\tan\theta$ 表示为
边 $DE$、$DF$、$EF$ 的比。
\end{exercise}

\begin{exercise}[$\star$]\label{exo:g9:trig:4}
在\omterm{def:g6:shapes:triangles}{直角三角形}中，斜边长为 $10$，一个锐角为
$28^\circ$。求两条直角边（\omterm{def:g6:decimals:rounding}{四舍五入}到十分位；
$\cos 28^\circ \approx 0.883$，$\sin 28^\circ \approx 0.469$）。
\end{exercise}

\begin{exercise}[$\star$]\label{exo:g9:trig:5}
在\omterm{def:g6:shapes:triangles}{直角三角形}中，角 $\theta$ 的邻边长为 $6$，
对边长为 $4.5$。求 $\tan\theta$，再求 $\theta$
（\omterm{def:g6:decimals:rounding}{四舍五入}到整度；$\tan^{-1}(0.75) \approx 36.9^\circ$）。
\end{exercise}

\begin{exercise}[$\star\star$]\label{exo:g9:trig:6}
一架无人机以 $120$ m 的高度飞行，地面观测者
看到它的仰角为 $32^\circ$。求无人机与观测者的
水平距离
（$\tan 32^\circ \approx 0.625$）。
\end{exercise}

\begin{exercise}[$\star\star$]\label{exo:g9:trig:7}
一条通道坡道须升高 $1.2$ m，且与水平面的夹角
至多为 $5^\circ$。求坡道沿斜面方向的最短长度
（$\sin 5^\circ \approx 0.0872$）。结果向上取到
最接近的十分之一米。
\end{exercise}

\begin{exercise}[$\star\star$]\label{exo:g9:trig:8}
利用\cref{prop:g9:trig:properties}：锐角 $\theta$ 满足
$\cos\theta = 0.6$。求 $\sin\theta$，不求 $\theta$，然后
求 $\tan\theta$。
\end{exercise}

\begin{exercise}[$\star\star\star$]\label{exo:g9:trig:9}
证明\cref{ex:g9:trig:special}中的精确值：在等腰\omterm{def:g6:shapes:triangles}{直角三角形}中，
若两直角边为 $1$，求斜边；在边长为 $1$ 的\omterm{def:g6:shapes:triangles}{等边三角形}中，
求高，并由此推出
$30^\circ$ 与 $60^\circ$ 的\omterm{def:g9:trig:ratios}{余弦}和\omterm{def:g9:trig:ratios}{正弦}。
\end{exercise}

\section{问题：测量无法到达之处}

\begin{problem}[{周末问题——测量员的双站法：测量不可达的高度、
道路坡度、到地平线的距离，
以及正切函数奔向无穷的竞赛}]
\label{pb:g9:trig:1}
卷尺够不到山峰的顶端、河对岸塔楼的顶部，
或海上的地平线——但测角仪与本章的三个比值可以，
本题的核心是测量员的经典方法：\emph{双站法}，
它不必接近底部便能算出高度；
周围还有路标、梯子、楼梯、地球的
曲率，以及初次见到一个函数
急剧增大至无穷。

\medskip
\noindent\textbf{第一部分——单站。}
（取值：$\tan 35^\circ \approx 0.700$，
$\sin 35^\circ \approx 0.574$，$\cos 35^\circ \approx 0.819$，
$\tan 1^\circ \approx 0.0175$。）
\begin{enumerate}
  \item 在距灯塔底部 $60$ m 处的一点，
        看到塔顶的仰角为 $35^\circ$（自水平方向、在眼睛高度处测量），
        求塔顶高出眼睛的
        高度？
  \item 一块路标标明坡度为 $12\,\%$：道路每升高 $12$ m，
        水平方向延伸 $100$ m。道路与水平面的夹角
        为多少？一条 $45^\circ$ 的道路，其坡度百\omterm{def:g9:fractions:fraction}{分数}
        是多少？
  \item \omterm{def:g6:shapes:triangles}{直角三角形}的两个锐角互余，且每个锐角的对边
        恰是另一个锐角的
        邻边。由此推出“余”恒等式：
        \[
        \sin(90^\circ - \theta) = \cos\theta,
        \qquad
        \cos(90^\circ - \theta) = \sin\theta .
        \]
  \item 一架 $8$ m 长的梯子与地面成 $60^\circ$ 的角斜靠在墙上。
        利用\cref{ex:g9:trig:special}中的精确值，
        求梯子能达到的高度，以及梯脚离墙的距离——
        先给出精确值，再精确到厘米。
  \item 验证 $\cos^2\theta + \sin^2\theta = 1$：对 $\theta = 35^\circ$
        作数值验证，对 $\theta = 60^\circ$ 作精确验证。
        为什么这一恒等式
        （见\cref{prop:g9:trig:properties}）可以
        作为检验计算器运算的好习惯？
\end{enumerate}

\medskip
\noindent\textbf{第二部分——双站。}
山顶可见，山脚无法到达。从
点 $A$ 测得山顶的仰角为 $30^\circ$；沿直线朝山走 $200$ m，
到达点 $B$，此时仰角变为
$40^\circ$。记 $h$ 为山顶高出眼睛的高度，
$x$ 为 $B$ 到山顶
铅垂线的水平距离。（取值：$\tan 30^\circ \approx 0.577$，
$\tan 40^\circ \approx 0.839$。）
\begin{enumerate}[resume]
  \item 将 $\tan 40^\circ$ 与 $\tan 30^\circ$ 分别用
        $h$、$x$ 与 $x + 200$ 表示：得到两个方程，两个未知数。
  \item 由每个方程用 $h$ 表示水平距离，
        两式相减，导出
        双站公式：
        \[
        h = \frac{200}
        {\ \dfrac{1}{\tan 30^\circ} - \dfrac{1}{\tan 40^\circ}\ } .
        \]
        求 $h$，精确到十米。
  \item 同时求 $x$，并将两个结果与
        $30^\circ$ 的视线（自 $A$ 点观测）相比较，做合理性检验。
  \item 用一句话说明：为什么双站法不可或缺——
        在此处无法测得的单个量，第 1 题的单站法
        本需要测得的是
        什么？
  \item 某塔的仰角先后测得为 $30^\circ$ 与 $60^\circ$，
        两次测量之间朝塔前进了 $200$ m。
        利用 $\tan 30^\circ = \frac{1}{\sqrt3}$ 与
        $\tan 60^\circ = \sqrt3$，求塔高的
        \emph{精确值}——答案是 $\sqrt3$ 的\omterm{def:g9:arith:divisor}{倍数}——
        再精确到米。
\end{enumerate}

\medskip
\noindent\textbf{第三部分——到地平线与沿墙向上。}
\begin{enumerate}[resume]
  \item 地平线有多远？站立时眼睛高于球形地球表面 $h$ 米，
        地球\omterm{def:g3:shapes:shapes}{半径}为 $R$。你的视线与球面相切，
        视线、通往切点的\omterm{def:g3:shapes:shapes}{半径}与
        脚下的\omterm{def:g3:shapes:shapes}{半径}构成一个\omterm{def:g6:shapes:triangles}{直角三角形}。
        利用毕达哥拉斯定理（见\cref{thm:g9:trig:pythagoras}）
        证明到地平线的距离 $d$
        满足 $d^2 = 2Rh + h^2 \approx 2Rh$，
        并计算 $d$，其中眼睛高度为 $1.80$ m
        （$R = 6.37 \times 10^6$ m）。
  \item 同样的问题，改在 $324$ m 高的塔顶求解。（水手的
        经验法则：以千米计的地平线距离约为
        $3.6\sqrt{h}$，$h$ 以米计——用它检验你的两个
        结果。）
  \item 拇指约宽 $2$ cm，伸直手臂举在离眼睛约 $60$ cm 处，
        它所张的角为多少？满月约张
        $0.5^\circ$，拇指中有多少
        能遮住整个满月？
  \item 舒适的楼梯踢面高 $17$ cm，踏面宽 $29$ cm。
        它们的爬升角为多少？将其与
        \cref{ex:g9:trig:special}中的 $30^\circ$ 相比，
        以及与倾角为 $60^\circ$ 的陡峭阁楼梯相比，
        \emph{它}在 $29$ cm 的踏面上需要怎样的踢面高度？
  \item \omterm{def:g9:trig:ratios}{正切}的竞赛：（用计算器）计算
        $\tan 10^\circ$、$\tan 45^\circ$、$\tan 80^\circ$、
        $\tan 89^\circ$。当角趋近于
        $90^\circ$ 时会发生什么？为何如此（想一想邻边）？墙面
        竖直，比值没有意义，图像急剧伸向
        无穷——你的第一条\emph{渐近线}，
        一种将在高中分册中细致研究的对象。
\end{enumerate}
\end{problem}
